\chapter{现代多重网格方法}
\label{ch:modern}

\begin{keybox}
\textbf{现代多重网格沿两条相互耦合的轴线发展：一条改变“怎样构造粗空间与校正”，另一条改变“怎样把这些操作映射到现代硬件”。}本章因此不按论文年份罗列，而是先补非线性、子空间与代数粗空间的数学扩展，再讨论低精度、专用矩阵单元、无矩阵高阶离散以及极大规模分布式粗网格算法。

阅读每一种方法时都回答四个问题：它放宽了前文哪一个假设？改变了粗空间、平滑器还是执行成本？引入了什么新的数值风险？应当用什么验证量确认收益？这样现代方法就能挂回前文已经建立的对象，而不会变成术语目录。带宽与算术强度见\chapref{ch:parallel-model}、\chapref{ch:mg-parallel-operations}和\chapref{ch:gpu}，AMG 层级与稀疏复杂度见\chapref{ch:software}，分布式粗层通信见\chapref{ch:mpi}和\chapref{ch:distributed-gpu}。
\end{keybox}

\section{现代多重网格的扩展方向}

前六章建立的经典路线可以概括为：线性离散系统、几何或代数层级、平滑器、粗网格校正以及递归 V-cycle。现代方法通常不是把这条链完全推翻，而是改变其中某一环：非线性问题需要改变“校正方程”的形式；复杂算子需要更自适应的粗空间；高阶离散改变算子的存储与应用方式；异构和极大规模机器则改变每一级操作的资源成本。

因此可以同时沿两类问题阅读本章。数学轴线问
\begin{equation}
\boxed{\text{什么误差由哪个子空间或粗空间表示，校正方程怎样定义？}}
\end{equation}
执行轴线问
\begin{equation}
\boxed{\text{这些操作怎样减少字节搬运、同步、通信或低并行度开销？}}
\end{equation}
两者不能完全分开：一个数学上更强的 smoother 可能具有更差的并行结构，一个更稀疏的粗算子也可能改变收敛因子。

\begin{warningbox}
本章只补足理解现代方法所需的最小模型，不试图替代非线性求解、有限元或 AMG 专著。凡是出现新的抽象，都要求读者至少能回答“未知量住在哪里、粗化改变什么、一次算子应用做什么”；论文级算法细节仍应回到专门文献。
\end{warningbox}

\section{非线性多重网格与 FAS}
\label{sec:modern-fas}

前文的 correction scheme 建立在线性系统
\begin{equation}
A_h u_h=f_h
\end{equation}
上。若当前近似为 $u_h$，残差 $r_h=f_h-A_hu_h$，粗层求的是误差近似 $e_H$，再 prolong 回细层。这一结构依赖于线性恒等式
\begin{equation}
A_h(u_h+e_h)=A_hu_h+A_he_h.
\end{equation}

对非线性离散方程
\begin{equation}
F_h(u_h)=f_h,
\label{eq:modern-nonlinear-problem}
\end{equation}
一般不存在可以独立于当前状态求解的线性误差方程。\textbf{全近似格式（Full Approximation Scheme, FAS）}改为在粗层直接求一个完整状态 $u_H$。设 $I_h^H$ 把细层状态限制到粗层，$R_h^H$ 把残差限制到粗层，则粗层右端写成
\begin{equation}
f_H^{FAS}
=
F_H(I_h^H u_h)
+
R_h^H\bigl(f_h-F_h(u_h)\bigr).
\label{eq:modern-fas-rhs}
\end{equation}
粗层求解
\begin{equation}
F_H(u_H)=f_H^{FAS},
\end{equation}
再用
\begin{equation}
u_h\leftarrow u_h+P_H^h\bigl(u_H-I_h^H u_h\bigr)
\label{eq:modern-fas-correction}
\end{equation}
更新细层。

这个写法最值得检查的是一个一致性极限：如果细层当前近似已经满足 $F_h(u_h)=f_h$，那么式~\eqref{eq:modern-fas-rhs}退化为
\begin{equation}
f_H^{FAS}=F_H(I_h^H u_h),
\end{equation}
于是 $u_H=I_h^H u_h$ 本身就是粗层解，式~\eqref{eq:modern-fas-correction}不给出额外修正。换言之，FAS 不会把一个已经满足细层方程的状态无故推离解。

在线性情形 $F_h(u)=A_hu$ 下，令
\begin{equation}
e_H=u_H-I_h^H u_h,
\end{equation}
代入粗层方程即可重新得到熟悉的 coarse-grid correction。因此 FAS 不是另一套无关算法，而是把前文线性 correction scheme 推广到“粗层也表示完整近似”的形式。

\section{子空间校正}
\label{sec:modern-subspace-correction}

前文把 smoother 和 coarse-grid correction 分开介绍，但两者可以放进同一个框架。设离散空间被若干子空间覆盖
\begin{equation}
V=V_0+V_1+\cdots+V_m,
\end{equation}
其中 $V_0$ 可以是粗空间，其他 $V_i$ 可以对应一个点、一条线、一个 patch 或一个子域。一次局部校正的本质是：\textbf{只在某个子空间内寻找能够降低当前误差的修正。}

点 Jacobi/Gauss--Seidel 可理解为极小局部子空间上的校正；line smoother 把一整条强耦合方向放进同一子空间；patch smoother 或 Schwarz 方法则在重叠局部区域上解小问题；coarse-grid correction 使用的是跨越全局低频误差的粗空间。这个视角解释了为什么“更强的局部求解”通常能处理更复杂的耦合，却也可能增加同步、局部因子分解和数据搬运成本。

若所有子空间校正并行计算后相加，得到的是 additive 结构；若前一个子空间的更新立即影响后一个子空间，则是 multiplicative 结构。现代并行实现经常在这两者之间做折中：数学上希望较强耦合，硬件上又希望减少顺序依赖。这与\chapref{ch:linear-solvers}中 Jacobi/Gauss--Seidel 的并行差异以及\chapref{ch:gpu}中的 patch/line kernel 设计是同一个矛盾在更一般空间中的表现。

\section{代数粗空间}
\label{sec:modern-amg-coarse-space}

几何多重网格利用坐标和规则粗化直接给出 $V_H$；AMG 则必须从离散算子本身判断哪些自由度应保留、哪些误差模式必须进入粗空间。经典 strength-of-connection 只是这一问题的一种答案。更一般地，现代 AMG 设计常围绕三个问题组织：
\begin{enumerate}
\item 当前 smoother 留下了哪些难以消除的误差？
\item 选出的 coarse variables / coarse basis 是否能表示这些误差？
\item prolongation 是否在保持近似能力的同时控制了算子复杂度和通信量？
\end{enumerate}

compatible relaxation、energy-minimizing interpolation、aggressive coarsening 等思想虽然具体算法不同，但都在回答上述问题。这里更重要的不是记住名称，而是把它们放回\chapref{ch:two-grid}的两条基本性质：\textbf{smoothing property 决定什么误差被留下，approximation property 决定粗空间能否表示这些误差。}现代 AMG 的许多“自适应”机制，本质上是在更复杂的系数、各向异性或图结构下重新估计这两个性质，而不是假定几何低频必然等于代数慢模态。

\section{混合精度多重网格}

对 memory-bound SpMV/stencil，把 FP64 数据换成 FP32，理论字节数近似减半；FP16 进一步下降。因此 mixed precision 的第一收益常来自 memory traffic，而不是低精度浮点峰值。

\paragraph{先做字节预算，再谈加速。}
若某向量场有 $N$ 个标量，仅比较存储本身，则 FP64、FP32、FP16 分别需要约 $8N$、$4N$、$2N$ byte。把一个层级的部分向量从 FP64 改成 FP32，最多只减少\textbf{这些被降精度对象}的字节数；index、halo metadata、未降精度系数场和通信同步并不会自动减半。因此“精度减半 $\Rightarrow$ solve time 减半”不是合法推论，必须回到\chapref{ch:benchmarking}测量完整 time-to-solution。

一个典型策略是：
\begin{equation}
\text{fine level: FP64/FP32},\qquad
\text{coarse levels: FP32/FP16},
\end{equation}
然后在 restriction/prolongation 时做 precision conversion。Tsai、Beams 与 Anzt 的 three-precision AMG 在 AMD、Intel、NVIDIA GPU 上研究了不同层级采用 double/single/half 的组合，说明 half precision 可以在保持最终精度的条件下用于部分 multigrid 层级\cite{tsai2023three}。

若把一个\textbf{固定线性、固定精度配置}的 V-cycle 近似看作平稳迭代，低精度执行可以示意为
\begin{equation}
\widehat E=E+\Delta E.
\end{equation}
在这种简化模型下，$\rho(\widehat E)<1$ 是误差传播保持收缩的自然判据。但这不是 mixed-precision multigrid 的普适定理：实际舍入误差依赖状态，scale/recover 操作可能非线性，inner tolerance 或 level precision 若随迭代变化，预条件器也可能成为 variable operator。

因此工程验证至少要同时看：
\begin{equation}
\boxed{
\text{residual history}
+
\text{iteration count}
+
\text{final high-precision residual/error}
+
\text{overflow/underflow behavior}.
}
\label{eq:modern-mixed-precision-validation}
\end{equation}
外层高精度 residual correction 或 flexible Krylov 可以容忍一定程度的低精度 inner solve，但前提应由实际收敛证据而不是单一谱半径近似判断。Krylov 子空间、CG/GMRES/FGMRES 与预条件的基础见\chapref{ch:linear-solvers}，它们与 multigrid 预条件器的组合见\chapref{ch:mg-cycles}的\secref{sec:mg-preconditioner}。

\subsection{FP16 结构化多重网格}

2024 年 Zong 等人的 structured multigrid preconditioner 工作并不是简单把整个 V-cycle 改成 FP16，而是针对 FP16 的数值范围和精度限制加入 scaling、setup-then-scale、运行时 recover/rescale，以及相应的数据格式/SIMD 实现\cite{zong2024fp16}。论文在其测试问题上报告了 preconditioner 与端到端流程的加速，但这些数字属于特定 CPU、矩阵和 workflow，不能直接作为 GPU 或其他应用的普适收益。

这项工作的可迁移机制是：低精度设计必须同时处理
\begin{equation}
\boxed{
\text{range}
+
\text{roundoff}
+
\text{conversion cost}
+
\text{memory format}.
}
\end{equation}

对工程实现，可把总误差做一个示意性预算：
\begin{equation}
\|e_{total}\|
\lesssim
C_d\|e_{disc}\|+
C_i\|e_{iter}\|+
C_r\|e_{round}\|.
\label{eq:modern-error-budget}
\end{equation}
这里 $C_d,C_i,C_r$ 依赖问题和所用范数，式~\eqref{eq:modern-error-budget}只是设计检查表，不是无条件误差定理。真正能否降低 coarse-level precision，应由式~\eqref{eq:modern-mixed-precision-validation}中的数值证据确认。

\section{Tensor Core 与代数多重网格}

\textbf{Tensor Core} 面向规则小块矩阵乘加；传统 CSR SpMV/SpGEMM 的不规则稀疏结构并不能直接高效映射过去。AmgT 的核心因此不是“把原来的 CSR kernel 换成 Tensor Core 指令”，而是先改变稀疏表示，再改变 setup/solve kernel\cite{lu2024amgt}。

作者公开的 AmgT 源码把实现直接嵌入一个 HYPRE fork\cite{amgtRepo2026}。仓库给出的三处关键路径是：
\begin{itemize}
\item \texttt{AmgT\_HYPRE/src/seq\_mv/seq\_mv.h}：mBSR 数据结构；
\item \texttt{AmgT\_HYPRE/src/seq\_mv/csr\_spgemm\_device.c}：setup 阶段 SpGEMM；
\item \texttt{AmgT\_HYPRE/src/seq\_mv/csr\_matvec\_device.c}：solve 阶段 SpMV。
\end{itemize}

因此源码层的心智模型是
\begin{equation}
\boxed{
\text{CSR/AMG hierarchy}
\rightarrow
\text{Tensor-Core-friendly sparse/block representation}
\rightarrow
\begin{cases}
\text{SpGEMM in setup},\\
\text{SpMV in solve}.
\end{cases}
}
\label{eq:modern-amgt-path}
\end{equation}

这项工作给出的可迁移问题不是“Tensor Core 是否比 CUDA core 快”，而是：
\begin{equation}
\boxed{
\text{格式转换 + padding + extra bytes}
<
\text{规则块计算带来的收益？}
}
\end{equation}
这个不等式必须把 setup 与 solve 一起测量；如果 hierarchy 只使用一次，格式转换成本尤其不能忽略。

\section{细粒度数据分块}

2024 年 SC Workshop 的 BrickLib GMG 工作在 NVIDIA、AMD、Intel GPU 系统上使用统一的 fine-grain brick data blocking，并把这种布局贯穿 stencil、inter-level transfer 与通信\cite{antepara2024brick}。论文在其三套测试系统上报告了 73\% 的 peak performance-portability metric，以及扩展到 512 GPUs 时约 87\% 的 weak-scaling parallel efficiency；这些数字是该实验矩阵中的结果，不应解释成 BrickLib 在任意机器上的固定效率。

更值得保留的是实现机制：BrickLib 不只优化单个 stencil 的局部访问顺序，还让 bricks 的物理排序服务通信，从而减少 on-node packing。也就是说，同一个 layout 同时影响
\begin{equation}
\boxed{
\text{stencil locality}
+
\text{transfer locality}
+
\text{MPI packing/communication}.
}
\label{eq:modern-brick-layout}
\end{equation}
这与\chapref{ch:architecture}的应用级所有权观点一致：数据布局若只服务计算 kernel，却让 phase boundary 付出昂贵转换，端到端性能仍可能变差。

\section{高阶离散与 $hp$ 粗化}
\label{sec:modern-high-order-minimal-model}

前文的七点 stencil 可以把每个未知量想成一个节点值或 cell 值；高阶有限元和 DG 的数据组织不同。一个网格单元 $K$ 内通常有多个局部自由度，它们描述某个次数为 $p$ 的多项式近似。算子应用不再只是读取固定几个几何邻居，而是对每个单元执行
\begin{equation}
\boxed{\text{取局部 DOF}\rightarrow\text{基函数/梯度运算}\rightarrow\text{局部积分}\rightarrow\text{散射回全局 DOF}.}
\end{equation}
连续有限元的相邻单元可能共享全局自由度；DG 则保留单元局部自由度，并通过数值通量耦合单元面。这里不需要掌握完整弱形式，只需要看到：\textbf{一个 element-local operator 往往可以由少量几何/系数和基函数规则现场重建，因此显式保存完整稀疏矩阵并不是唯一选择。}

高阶问题还引入了两个不同的粗化方向：
\begin{itemize}
\item $h$-coarsening：减少网格单元数，几何尺度变粗；
\item $p$-coarsening：保留网格拓扑，但降低每个单元的多项式次数和局部自由度数。
\end{itemize}
例如从 $p=7$ 降到 $p=3$，几何单元没有消失，但每个单元的近似空间显著缩小。这就是后文 $hp$-multigrid 中字母 $h$ 与 $p$ 分别代表的结构变化。

\begin{figure}[htbp]
\centering
\begin{tikzpicture}[scale=0.9]
  \draw[thick] (0,0) rectangle (3,3);
  \draw (1.5,0)--(1.5,3); \draw (0,1.5)--(3,1.5);
  \foreach \x in {0.3,0.9,1.8,2.4} {
    \foreach \y in {0.3,0.9,1.8,2.4} {\fill (\x,\y) circle (1.2pt);}
  }
  \node[below] at (1.5,-0.25) {$h$-粗化：减少单元};
  \draw[->,thick] (3.5,1.5)--(5,1.5);
  \draw[thick] (5.5,0.75) rectangle (7,2.25);
  \foreach \x in {5.8,6.7} {\foreach \y in {1.05,1.95}{\fill (\x,\y) circle (1.2pt);}}

  \draw[thick] (9,0) rectangle (12,3);
  \foreach \x in {9.3,9.9,10.5,11.1,11.7}{\foreach \y in {0.3,0.9,1.5,2.1,2.7}{\fill (\x,\y) circle (1.0pt);}}
  \draw[->,thick] (12.5,1.5)--(14,1.5);
  \draw[thick] (14.5,0) rectangle (17.5,3);
  \foreach \x in {15.1,16.0,16.9}{\foreach \y in {0.6,1.5,2.4}{\fill (\x,\y) circle (1.2pt);}}
  \node[below] at (13.25,-0.25) {$p$-粗化：同一单元上减少局部 DOF};
\end{tikzpicture}
\caption{$h$-粗化与 $p$-粗化改变的是不同结构。示意图只表达“单元数量”与“每单元近似空间维数”的区别，不代表具体有限元节点布局。}
\label{fig:modern-h-p-coarsening}
\end{figure}

sum factorization 的作用可以在这个最小模型上理解：若高维基函数具有张量积结构，就不必把一个 $d$ 维局部算子当成稠密大矩阵直接乘，而可以分解成若干一维变换，从而减少计算复杂度与数据搬运。它并不改变离散问题，只改变 element-local operator 的执行方式。

\section{无矩阵高阶多重网格}

在\secref{sec:modern-high-order-minimal-model}的高阶离散中，显式矩阵会产生大量内存流量，无矩阵算子可以用更多局部 FLOP 换更少 bytes。Cui 与 Kanschat 的 2025 年 GPU Stokes multigrid 是一个具体实现例子：离散采用 $H(\mathrm{div})$-conforming DG，smoother 是 multiplicative Schwarz 的 vertex-patch 局部问题，局部求解结合 Schur complement 与 fast diagonalization，operator apply 则采用无矩阵计算与 conflict-free shared-memory access；论文还测试 mixed precision\cite{stokesgpu2025}。

这里可以直接看到数学结构与 GPU kernel 的对应：
\begin{equation}
\boxed{
\text{vertex-patch subspace}
\rightarrow
\text{local tensor-product solve}
\rightarrow
\text{fast diagonalization}
\rightarrow
\text{shared-memory kernel}.
}
\label{eq:modern-stokes-source-model}
\end{equation}

对于低阶七点 Poisson，增加局部 FLOP 的收益空间通常更小；对高阶 tensor-product operator，sum factorization、fast diagonalization 和 element-local reconstruction 则可能显著提高 arithmetic intensity。是否真正跨过 Roofline ridge point，仍需测实际 bytes 与 FLOPs。

\subsection{无矩阵代数 $hp$-多重网格}

2026 年 Ohm、Harper 与 Jansson 在 Neko 中研究了无矩阵代数 $hp$-多重网格：先通过 $p$-coarsening 从高阶谱元离散降到低阶，再仅利用 mesh adjacency 构造空间 aggregation，从而不依赖显式矩阵条目或几何 rediscretization\cite{ohm2026hpmg}。

论文实现给出了比“使用 Chebyshev smoother”更具体的层级策略：$p$-multigrid levels 使用二阶 Jacobi-accelerated Chebyshev，而 algebraic levels 使用三阶 Chebyshev；其目的之一是避免在完全无矩阵的 AMG levels 保存额外局部矩阵/对角信息。它还用 power iteration 估计 Chebyshev 所需谱区间。这里的关键不是这些阶数本身，而是\textbf{smoother 选择受到可获得 operator metadata 的直接约束}。

这项工作说明了一种可迁移设计：
\begin{equation}
\boxed{
\text{高阶无矩阵}
\rightarrow
p\text{-coarsening}
\rightarrow
\text{low-order adjacency}
\rightarrow
\text{无矩阵代数 }h\text{-粗化}.
}
\label{eq:modern-hpmg-path}
\end{equation}
它不能单独证明“Chebyshev 总比 Gauss--Seidel 更适合 GPU”，但确实展示了在无矩阵 hierarchy 中，规则多项式 smoother 可以避免依赖显式 sparse matrix entries。

\section{Exascale 粗网格通信}

\textbf{Exascale} 通常指面向 $10^{18}$ FLOP/s 量级超级计算系统的算法与软件环境。在这一尺度上，单个 GPU 越快，网络延迟和全局同步的相对成本越显眼。多重网格优化很大一部分并不是继续优化 fine-grid stencil，而是：
\begin{itemize}
\item aggressive process/rank coarsening；
\item topology-aware redistribution；
\item non-Galerkin sparsification，控制 coarse operator complexity；
\item communication-avoiding / asynchronous smoothers；
\item 减少 global reductions；
\item 让 NIC、GPU 和 MPI runtime 更好地异步协同。
\end{itemize}

这也解释了为什么“单 GPU kernel 已达到很高带宽”并不代表 distributed MLMG 已经解决。

\section{本章小结}
现代多重网格可以沿“数学空间”和“执行表示”两条轴理解。FAS 改变非线性粗层方程，子空间校正统一 point/line/patch/coarse correction，现代 AMG 与 $hp$ 方法重新设计粗空间；mixed precision、AmgT、BrickLib、高阶无矩阵 kernels 与 distributed coarse strategies 则改变数据表示和机器成本。

本章更重要的阅读方法是区分三层证据：\textbf{论文在特定实验上的结果}、\textbf{源码中真正采用的数据结构/算法路径}、以及\textbf{可以迁移到其他求解器的机制}。例如 AmgT 可迁移的是“先重构 sparse representation 再利用矩阵专用硬件”，而不是某个固定加速比；Neko hpMG 可迁移的是“无矩阵表示可提供的 metadata 会反过来约束 smoother 与 coarse-space construction”，而不是某个固定 Chebyshev 阶数。

\section*{习题}

\subsection*{现代方法的数学与性能模型}
\begin{enumerate}[label=\arabic*.]
\item 一个三维高阶单元在每个方向使用 $p+1$ 个局部基函数。分别计算 $p=7$ 与 $p=3$ 时每单元局部 DOF 数，并估计一次 $p$-coarsening 后局部未知量减少比例。
\item 一个层级数据集包含 $4N$ 个 FP64 标量。计算其存储量；若其中 $2N$ 个改为 FP32，计算新存储量和理论 bytes 减少比例。
\item 一个 memory-bound V-cycle 每次从 DRAM 搬运 800 GB 数据，硬件可持续带宽 1.6 TB/s。若 mixed precision 将流量降到 520 GB，但增加 15\% FLOP 且仍 bandwidth-bound，估算理想时间变化。
\item 某 low-precision coarse correction 每次把 residual 降低到原来的 0.25，但随后高精度 refinement 有固定 0.1 ms 成本。计算达到 $10^{-8}$ 需要多少 correction steps，并建立总时间表达式。
\item 高阶无矩阵 Kernel arithmetic intensity 为 12 FLOP/byte，低阶 CSR 为 0.3 FLOP/byte。对 ridge point 8 FLOP/byte 的 GPU 判断二者分别处在哪个 Roofline 区域。

\item 从式~\eqref{eq:modern-fas-rhs}出发，在线性情形 $F_h(u)=A_hu$ 下令 $e_H=u_H-I_h^H u_h$，推导它如何退化为标准 coarse-grid correction，并指出需要怎样选择 coarse operator 才与前文 Galerkin 结构一致。
\item 对一个由点子空间、line 子空间和 coarse space 组成的抽象分解，解释 additive 与 multiplicative 校正的执行依赖差异，并判断哪一种更容易并行。
\item 若低精度 correction 满足 $\|e_{k+1}\|\le q\|e_k\|$、$q<1$，推导达到给定误差阈值所需 correction 次数，并说明何时高精度 outer refinement 可以容忍低精度 inner solve。
\item 从 bytes/FLOP 模型说明 mixed precision 对 memory-bound multigrid 首先影响的是数据流量，而不是“自动增加 FLOP/s”。
\item 解释 Tensor Core/矩阵专用硬件为什么不能直接加速任意稀疏 SpMV；从数据布局与运算形态给出必要条件。
\item 对高阶无矩阵，说明“增加 FLOP、减少 bytes”为什么可能提高 GPU 利用率；用 Roofline 公式给出条件。
\end{enumerate}

\subsection*{数值与性能实验}
\begin{enumerate}[label=\arabic*.]
\item 将 V-cycle 的部分 level/向量改为 FP32，其余保持 FP64，比较 residual history、最终误差、bytes 与 wall time。
\item 实现 mixed-precision iterative refinement：低精度求 correction，高精度计算 residual。扫描 inner tolerance，观察稳定区间。
\item 对同一高阶/低阶算子比较 explicit sparse 与无矩阵apply 的 arithmetic intensity 和实测性能。
\item 若硬件支持，实验一个 block-dense/tensor-core-friendly smoother，与普通 sparse smoother 比较 setup 与 solve 成本。
\end{enumerate}

\subsection*{研究比较与技术选择}
\begin{enumerate}[label=\arabic*.]
\item mixed precision 节省 35\% 时间但使迭代次数增加 20\%。说明还需要哪些 accuracy/robustness 证据才能决定是否采用。
\item GPU 算力继续提高但网络 latency 基本不变时，解释为什么 coarse-grid communication 会成为更显著的瓶颈。
\item 选择无矩阵、mixed precision、communication-avoiding 中任两种现代技术，分析它们可能互补还是互相冲突，并设计联合 benchmark。
\end{enumerate}
